\documentclass[11pt]{article}

\usepackage[margin=1in]{geometry}
\usepackage{amsmath,amssymb,amsthm,mathtools}
\usepackage{graphicx}
\usepackage{booktabs}
\usepackage{enumitem}
\usepackage[hypertexnames=false]{hyperref}
\usepackage{url}
\usepackage[T1]{fontenc}
\usepackage[utf8]{inputenc}
\usepackage{lmodern}
\usepackage[backend=biber,style=numeric,sorting=nyt]{biblatex}
\addbibresource{research_guide_references.bib}
% Ensure that 'research_guide_references.bib' exists in the same directory as this .tex file.
% If it is located elsewhere, update the path accordingly, e.g.:
% \addbibresource{../references/research_guide_references.bib}

\title{Research Guide: Graph-Smooth Sparse Orthogonal PCA Beyond Graph-Regularized PCA}
\author{Internal Research Document}
\date{\today}

\begin{document}

\maketitle

\section{Executive Summary}

Recent work on Graph-Regularized PCA (GR-PCA) provides a principled framework for incorporating feature dependency structure into PCA through a learned feature graph and Laplacian regularization \cite{briola2026grpca}. The method is an important contribution to graph-aware dimensionality reduction, especially in settings where graph-coherent low-frequency structure is more meaningful than purely variance-maximizing directions. However, GR-PCA does not by itself solve the more specific problem of \emph{sparse, simultaneously orthogonal, graph-smooth multi-component PCA} under a modern manifold-based optimization framework \cite{chen2020amanpg,chen2024nonsmoothstiefel}.

This document presents:
\begin{itemize}
    \item a revised gap analysis of GR-PCA and related graph-aware PCA methods,
    \item a clearer separation between what is already known and what remains underexplored,
    \item a proposed research direction centered on \textbf{graph-smooth sparse orthogonal PCA} rather than a broad claim about graph-aware sparse PCA in general.
\end{itemize}

\section{Problem Setting}

Given centered data matrix $X \in \mathbb{R}^{n \times p}$ with sample covariance
\[
\widehat{\Sigma} = \frac{1}{n} X^\top X,
\]
classical PCA solves
\[
\max_{\|v\|_2 = 1} v^\top \widehat{\Sigma} v.
\]

Sparse PCA introduces sparsity, for example through an $\ell_1$ penalty \cite{zou2006spca}:
\[
\max_{\|v\|_2 = 1} v^\top \widehat{\Sigma} v - \lambda \|v\|_1.
\]

Graph-aware PCA augments this with structural regularization derived from a feature graph. In GR-PCA, the formulation is written in low-rank reconstruction form as \cite{briola2026grpca}
\[
\min_{U,V} \frac{1}{2}\|X - UV^\top\|_F^2
+ \alpha \sum_{j=1}^p \frac{\|V_{j:}\|_1}{1+d_j}
+ \frac{\lambda}{2}\operatorname{tr}(V^\top L V),
\]
where $L$ is the Laplacian of a graph inferred from a sparse precision matrix.

\section{What GR-PCA Already Solves Well}

GR-PCA has several genuine strengths \cite{briola2026grpca}:
\begin{itemize}
    \item it incorporates feature dependency structure through a graph learned from the data,
    \item it combines sparsification and Laplacian regularization in one objective,
    \item it is well motivated in regimes with graph-correlated high-frequency noise,
    \item it produces graph-coherent and interpretable loading patterns in synthetic experiments.
\end{itemize}

Accordingly, the goal of the present research should \emph{not} be framed as replacing GR-PCA or claiming that graph-aware sparse PCA is absent from the literature. The more appropriate goal is to study a narrower and still underdeveloped intersection: sparse multi-component loading estimation that is simultaneously graph-smooth and explicitly orthogonal.

\section{Revised Critical Gaps}

\subsection{Gap 1: No Explicit Multi-Component Orthogonality Constraint}

GR-PCA does not explicitly impose
\[
V^\top V = I.
\]

This matters because:
\begin{itemize}
    \item orthogonality is central to the classical interpretation of PCA,
    \item without explicit orthogonality, multiple components may be rotation-ambiguous,
    \item graph-coherent sparsification and multi-component decorrelation are left to interact implicitly rather than by design.
\end{itemize}

\textbf{Revised opportunity:} develop a formulation that explicitly enforces simultaneous orthogonality while retaining graph-aware sparse estimation \cite{chen2020amanpg,bertsimas2023geospca,cheng2026orthogonalspca}.

\subsection{Gap 2: Two-Stage Graph Construction}

GR-PCA first estimates a sparse precision matrix, then thresholds or extracts its support to form a feature graph, and only afterward solves the graph-regularized PCA problem \cite{briola2026grpca}.

Pipeline:
\begin{enumerate}
    \item estimate sparse precision matrix $\widehat{\Theta}$,
    \item construct graph from its support,
    \item perform graph-regularized PCA using the resulting Laplacian.
\end{enumerate}

This creates several limitations:
\begin{itemize}
    \item graph estimation error can propagate into the PCA stage,
    \item uncertainty in graph recovery is not modeled jointly with component estimation,
    \item thresholding may discard useful continuous edge-strength information.
\end{itemize}

\textbf{Revised opportunity:} study joint, alternating, or bilevel graph-plus-component estimation, or explicitly analyze robustness to graph misspecification.

\subsection{Gap 3: The Graph Prior Is Primarily a Smoothness Prior}

The core structural regularizer in GR-PCA is
\[
\operatorname{tr}(V^\top L V),
\]
which is a Laplacian smoothness penalty \cite{briola2026grpca}.

This is useful, but it encodes a specific inductive bias:
\begin{itemize}
    \item it favors low-frequency, graph-smooth loading patterns,
    \item it does not explicitly model signed antagonistic structure,
    \item it does not directly enforce localized discontinuities, subgraph selection, or connected-support constraints.
\end{itemize}

Thus, the limitation is \emph{not} that the method ignores graph structure, but that it relies on one particular notion of graph structure: smoothness.

\textbf{Revised opportunity:}
\begin{itemize}
    \item signed or weighted Laplacian variants,
    \item graph total variation penalties,
    \item connected-support or subgraph-aware regularizers.
\end{itemize}

\subsection{Gap 4: Degree-Weighted Sparsity Is Not the Same as Structured Support Recovery}

GR-PCA uses the weighted sparsity term
\[
\sum_{j=1}^p \frac{\|V_{j:}\|_1}{1+d_j},
\]
which penalizes low-degree features more strongly than high-degree ones \cite{briola2026grpca}.

This is a deliberate design choice, but it has consequences:
\begin{itemize}
    \item it biases sparsity according to node degree,
    \item it does not explicitly enforce connected or subgraph-structured supports,
    \item it may be undesirable in applications where degree should not influence feature selection strength.
\end{itemize}

\textbf{Revised opportunity:}
\begin{itemize}
    \item graph-structured sparsity penalties,
    \item overlapping-group penalties,
    \item connected-support constraints,
    \item row-sparse or component-shared support models when multiple components are expected to share variables.
\end{itemize}

\subsection{Gap 5: Evaluation Is Primarily Synthetic}

The current GR-PCA study is built around synthetic experiments across graph topologies, noise regimes, and sparsity levels \cite{briola2026grpca}.

This is valuable for isolating mechanisms, but it leaves open:
\begin{itemize}
    \item whether the method remains effective on noisy real feature graphs,
    \item whether the graph prior is robust outside the controlled simulation setting,
    \item whether interpretability gains persist in real scientific applications.
\end{itemize}

\textbf{Revised opportunity:}
\begin{itemize}
    \item gene expression with pathway or PPI graphs,
    \item neuroimaging with spatial or connectivity graphs,
    \item financial data with sector or supply-chain graphs,
    \item sensor networks with spatial adjacency.
\end{itemize}

\subsection{Gap 6: Baseline Coverage Is Limited}

GR-PCA is compared primarily against standard PCA and SparsePCA-style baselines \cite{briola2026grpca}.

This leaves out several relevant comparison classes:
\begin{itemize}
    \item graph + sparsity PCA variants \cite{feng2019glspca,wu2019gdspca,wang2025rsspca},
    \item joint orthogonal sparse PCA methods \cite{bertsimas2023geospca,cheng2026orthogonalspca},
    \item manifold-based sparse PCA solvers \cite{chen2020amanpg},
    \item structured sparse PCA baselines \cite{jenatton2010sspca}.
\end{itemize}

\textbf{Revised opportunity:} benchmark against a stricter baseline set so the contribution is evaluated relative to the most relevant parts of the literature rather than only to unstructured PCA.

\subsection{Gap 7: Theory Is Incomplete Rather Than Absent}

GR-PCA has a clear structural modeling rationale, but several questions remain open:
\begin{itemize}
    \item optimization convergence guarantees,
    \item statistical consistency and support recovery theory,
    \item robustness under graph misspecification,
    \item sensitivity to graph estimation error.
\end{itemize}

This should be framed as \emph{incomplete theory}, not as a total lack of theory.

\textbf{Revised opportunity:} provide convergence-to-stationarity guarantees for a proposed algorithm and, if possible, perturbation or robustness analysis with respect to graph errors.

\subsection{Gap 8: The Optimization Challenge Is Stronger Than ``Gradient + Projection''}

A key methodological caution is that the desired extension
\[
\max_{V^\top V = I}
\operatorname{tr}(V^\top \widehat{\Sigma} V)
- \lambda_1 \|V\|_1
- \lambda_2 \operatorname{tr}(V^\top L V)
\]
is not naturally solved by simply applying Euclidean soft-thresholding and then projecting onto the Stiefel manifold.

Why?
\begin{itemize}
    \item orthogonality introduces nonconvex manifold geometry,
    \item the Laplacian term couples coefficients across features,
    \item projection or retraction after naive thresholding may distort sparsity patterns.
\end{itemize}

\textbf{Revised opportunity:} formulate the method using a split-variable or manifold proximal framework rather than an overly naive projected-gradient description \cite{chen2020amanpg,chen2024nonsmoothstiefel,zhu2025stccal,li2024structurednonsmooth}.

\section{Unified Open Problem}

The meaningful open problem is not
\begin{quote}
``make graph-aware PCA sparse and gradient-based,''
\end{quote}
because partial versions of that already exist \cite{feng2019glspca,wu2019gdspca,wang2025rsspca}.

A better formulation is:

\begin{quote}
Design a PCA-specific method that is simultaneously
\begin{itemize}
    \item sparse,
    \item smooth on a known feature graph,
    \item explicitly orthogonal across multiple components,
    \item computationally practical,
    \item and supported by convergence and robustness analysis.
\end{itemize}
\end{quote}

\section{Revised Proposed Research Direction}

\subsection{Safer Core Formulation}

A more optimization-friendly formulation is a split-variable model:
\[
\min_{A^\top A = I_r,\; B}
-\operatorname{tr}(A^\top \widehat{\Sigma} A)
+ \lambda_1 \|B\|_1
+ \lambda_2 \operatorname{tr}(B^\top L B)
+ \frac{\rho}{2}\|A-B\|_F^2.
\]

Interpretation:
\begin{itemize}
    \item $A$ carries orthogonality on the Stiefel manifold,
    \item $B$ carries sparse graph-regularized structure in Euclidean space,
    \item the quadratic coupling links the two.
\end{itemize}

This avoids presenting the problem as a trivial projected proximal-gradient extension.

\subsection{Optimization Approach}

A more defensible solver description is:
\begin{itemize}
    \item manifold-based updates for the orthogonality variable,
    \item graph-coupled sparse updates for the Euclidean variable,
    \item alternating optimization with convergence-to-stationarity as the initial theoretical goal.
\end{itemize}

Possible solver choices include:
\begin{itemize}
    \item first-order proximal or primal-dual updates for the graph-coupled sparse step,
    \item semismooth Newton or related acceleration only as an extension,
    \item retraction-based manifold updates for the orthogonality step.
\end{itemize}

\subsection{Revised Contribution Targets}

The most realistic contribution targets are:
\begin{enumerate}
    \item a PCA-specific formulation combining graph smoothness, sparsity, and explicit orthogonality,
    \item a practical alternating solver,
    \item a rigorous empirical study of the variance--structure--interpretability tradeoff,
    \item supportive convergence analysis to stationary points.
\end{enumerate}

The paper should \emph{not} be positioned as the first instance of graph-coupled nonsmooth manifold optimization in the abstract, because adjacent work already combines several of these ingredients outside PCA \cite{zhu2025stccal}.

\section{Revised Theoretical Roadmap}

A realistic theorem roadmap is:
\begin{itemize}
    \item lower boundedness and bounded iterates,
    \item sufficient decrease of the alternating objective,
    \item vanishing successive differences,
    \item stationarity of accumulation points,
    \item optional KL-style full-sequence convergence if cleanly provable.
\end{itemize}

Avoid leading with:
\begin{itemize}
    \item global optimality,
    \item broad new manifold theory,
    \item universal superiority over prior methods.
\end{itemize}

\section{Revised Experimental Roadmap}

\subsection{Baselines}

A stronger baseline set should include:
\begin{itemize}
    \item PCA,
    \item standard sparse PCA,
    \item GR-PCA,
    \item graph + sparse PCA baselines,
    \item joint orthogonal sparse PCA baselines,
    \item your own ablations:
    \begin{itemize}
        \item no graph term,
        \item no sparsity term,
        \item deflation instead of joint extraction.
    \end{itemize}
\end{itemize}

\subsection{Metrics}

Recommended metrics:
\begin{itemize}
    \item explained variance,
    \item reconstruction error,
    \item support recovery (precision, recall, F1),
    \item subspace distance on synthetic data,
    \item graph smoothness / Laplacian energy,
    \item support connectivity metrics,
    \item runtime and memory,
    \item sensitivity to graph corruption.
\end{itemize}

\subsection{Datasets}

Use:
\begin{itemize}
    \item synthetic graph-aligned data for controlled recovery studies,
    \item one or two real domains only, chosen to match the graph prior.
\end{itemize}

\section{Revised Positioning Statement}

\textbf{GR-PCA is a strong graph-aware PCA baseline, but it does not directly solve the specific problem of sparse, simultaneously orthogonal, graph-smooth multi-component PCA within a modern manifold-based optimization framework. Its graph prior is centered on Laplacian smoothness, its sparsity is degree-weighted rather than explicitly connected-support structured, and its empirical study is still primarily synthetic with limited baseline coverage.}

This work should therefore be positioned as a targeted extension at the intersection of:
\begin{itemize}
    \item sparse PCA,
    \item graph-aware PCA,
    \item joint orthogonal component extraction,
    \item and structured manifold optimization,
\end{itemize}
rather than as a broad claim that these ingredients have never been combined before.

\section{Conclusion}

Graph-aware PCA is no longer an empty research area, and graph+sparsity PCA is not absent from the literature. The defensible opportunity is narrower and more interesting: to study how sparse, explicitly orthogonal, graph-smooth multi-component PCA can be formulated and optimized in a way that is both practically useful and theoretically well supported. GR-PCA is best viewed as a strong foundation and comparison point, not as a strawman.

\printbibliography

\end{document}